\documentclass[12pt]{article}
\usepackage[a4paper,margin=1in]{geometry}
\usepackage{amsmath}
\usepackage{booktabs}

\begin{document}
	
	\begin{center}
		\Large \textbf{Determination of the Composition of an Unknown Binary Liquid Mixture}\\
	\end{center}
	
	\section*{Aim}
	
	To determine the composition of an unknown mixture of two liquids by measuring its refractive index using an Abbe refractometer.
	
	\section*{Theory}
	
	The refractive index of a liquid is defined as the ratio of the velocity of light in vacuum to its velocity in the liquid:
	
	\[
	n = \frac{c}{v}
	\]
	
	where \(c\) is the velocity of light in vacuum and \(v\) is the velocity of light in the liquid.
	
	The refractive index of a binary liquid mixture depends on its composition, provided that the temperature and wavelength of light remain constant. Therefore, the refractive indices of mixtures having known compositions can be measured and plotted as a calibration graph. \\	
	The refractive index of the unknown mixture is measured under the same experimental conditions. Its composition is found from the calibration graph.
	
	\section*{Apparatus Required}
	
	Abbe refractometer, lens tissue, dropper, standard mixtures of liquids A and B, and unknown liquid mixture.
	
	\section*{Procedure}
	
	\begin{enumerate}
		\item Clean the upper and lower prism surfaces of the Abbe refractometer using lens tissue and a suitable solvent.
		
		\item Maintain a constant temperature throughout the experiment and record the temperature.
		
		\item Prepare standard mixtures of liquids A and B with known compositions, such as 0, 10, 20, 30, 40, 50, 60, 70, 80, 90, and 100\% of liquid A by volume.
		
		\item Place 2--3 drops of the first standard mixture on the lower prism and close the upper prism carefully. Ensure that a uniform liquid film is formed without air bubbles.
		
		\item Illuminate the prism and observe through the eyepiece.
		
		\item Adjust the compensator until the light-dark boundary becomes sharp and colourless.
		
		\item Adjust the refractometer knob until the light-dark boundary coincides with the crosshair.
		
		\item Record the refractive index. Take three readings and calculate the mean refractive index.
		
		\item Clean the prism surfaces and repeat the procedure for all standard mixtures.
		
		\item Plot a calibration graph with percentage composition of liquid A on the \(x\)-axis and mean refractive index on the \(y\)-axis.
		
		\item Place the unknown mixture on the prism and measure its refractive index.
		
		\item Locate the refractive index of the unknown on the \(y\)-axis of the calibration graph. Draw a horizontal line to meet the calibration curve and then a vertical line to the \(x\)-axis.
		
		\item Read the percentage of liquid A from the graph. The percentage of liquid B is calculated using:
		\[
		\%B = 100 - \%A
		\]
	\end{enumerate}
	
	\section*{Observation Table}
	
	\begin{center}
		\begin{tabular}{cccc}
			\toprule
			Mixture & Liquid A (\% v/v) & Liquid B (\% v/v) &
			Refractive Index, \(n_D\) \\
			\midrule
			1 & 0   & 100 & \\
			2 & 10   & 90 & \\
			3 & 20   & 80 & \\
			4 & 30   & 70 & \\
			5 & 40   & 60 & \\
			6 & 50  & 50  & \\
			7 & 60  & 40  & \\
			8 & 70  & 30  & \\
			9 & 80  & 20  & \\
			10 & 90   & 10 & \\
			11 & 100 & 0   & \\
			Unknown & -- & -- & \\
			\bottomrule
		\end{tabular}
	\end{center}
	
	\section*{Result}
	
	From the calibration graph:
	
	\[
	\text{Composition of liquid A} = \underline{\hspace{2cm}} \% \text{ (v/v)}
	\]
	
	\[
	\text{Composition of liquid B} = \underline{\hspace{2cm}} \% \text{ (v/v)}
	\]
	
	\section*{Precautions}
	
	\begin{itemize}
		\item Keep the temperature constant for all measurements.
		\item Clean the prism surfaces before measuring each sample.
		\item Avoid formation of air bubbles between the prism surfaces.
		\item Ensure that the light-dark boundary is sharp and colourless before taking readings.
		\item Use identical experimental conditions for standard and unknown mixtures.
	\end{itemize}
	
\end{document}